\chapter{Despliegue en servidor}\label{cap:despliegue}

Este capítulo resume el despliegue por versiones: organización de la
instalación, servicios, construcción del paquete, activación y reversión. Es
un procedimiento genérico; los valores propios de cada servidor aparecen
como marcadores en mayúsculas. La guía completa es
\archivo{docs/deployment.md}.

\section{Modelo de despliegue}\label{sec:desp-modelo}

Cada versión es un paquete inmutable identificado por etiqueta y commit,
sin credenciales, datos ni claves. El estado persistente vive fuera de las
versiones: cambiar de versión equivale a mover un enlace simbólico y
reiniciar dos servicios. Todo se ejecuta con una cuenta Unix propia, sin
privilegios de administrador, bajo un directorio raíz privado (en adelante
\texttt{RAIZ}).

\begin{table}[H]
  \centering
  \caption{Estructura del directorio raíz del despliegue.}\label{tab:desp-raiz}
  \begin{tabularx}{\textwidth}{L{5.1cm} Y}
    \toprule
    \textbf{Ruta bajo RAIZ} & \textbf{Contenido} \\
    \midrule
    \texttt{config/}
      & \texttt{settings.json}, \texttt{nginx.conf}, \texttt{redis.conf} y
        la huella de migraciones admitida (\texttt{schema}). Véase el
        \cref{cap:configuracion}. \\
    \texttt{data/}
      & Bases de datos, material de la CA y de la TSA, y temporales. \\
    \texttt{releases/}
      & Un directorio \texttt{VERSION-COMMIT} por versión admitida, con su
        manifiesto \texttt{release.json}. \\
    \texttt{current}, \texttt{previous}
      & Enlaces simbólicos a la versión activa y a la anterior. \\
    \texttt{incoming/}
      & Paquetes recibidos, pendientes de admisión. \\
    \texttt{tools/}
      & Controladores Python instalados (copia de \texttt{ops/deploy}). \\
    \texttt{backups/}, \texttt{logs/}, \texttt{run/}
      & Respaldos, registros de nginx y archivos de ejecución. \\
    \texttt{maintenance/}
      & Diarios de las operaciones sobre los controladores. \\
    \texttt{controller\_launcher.py}
      & Lanzador de controladores, fuera de \texttt{tools/}. \\
    \texttt{deploy.lock}
      & Bloqueo exclusivo: una sola operación de despliegue a la vez. \\
    \bottomrule
  \end{tabularx}
\end{table}

\section{Servicios y proxy inverso}\label{sec:desp-servicios}

La aplicación son cuatro unidades systemd de usuario que escuchan solo en
\texttt{127.0.0.1}. Los puertos de la \cref{tab:desp-unidades} son los que
\archivo{ops/deploy/host.py} escribe por omisión en \texttt{settings.json}.

\begin{table}[H]
  \centering
  \caption{Unidades de servicio de la aplicación.}\label{tab:desp-unidades}
  \begin{tabularx}{\textwidth}{L{4.1cm} L{1.5cm} Y}
    \toprule
    \textbf{Unidad} & \textbf{Puerto} & \textbf{Función y dependencias} \\
    \midrule
    \texttt{qadra-postgres.service} & 15486
      & PostgreSQL propio de la instalación. \\
    \texttt{qadra-redis.service} & 16386
      & Redis propio, con contraseña y persistencia AOF. \\
    \texttt{qadra-api.service} & 18087
      & API del sistema. Requiere las dos bases. \\
    \texttt{qadra-web.service} & 18086
      & nginx: frontend y proxy hacia la API, a la que espera. \\
    \bottomrule
  \end{tabularx}
\end{table}

El proxy inverso es el nginx de \texttt{qadra-web}
(\texttt{RAIZ/config/nginx.conf}): sirve \texttt{current/web}, reenvía
\texttt{/api/} a la API y responde 403 a \texttt{/api/v1/auth/bootstrap},
de modo que el alta del primer Administrador solo puede hacerse en el
propio servidor, contra el puerto de la API. La unidad \texttt{qadra-api} no ejecuta el
binario directamente: lanza, por medio del lanzador en las unidades
aprobadas, el controlador \texttt{tools/runtime.py}, que arranca \texttt{despacho-cli serve} con las rutas y los secretos de la
instalación. El portal de la CA tiene unidades propias (\cref{cap:ca}).

\nota{Publicar la aplicación por
HTTPS (dominio, certificado TLS y proxy público hacia el puerto de
\texttt{qadra-web}) es configuración del servidor ajena a estos guiones; sin
ella se accede por un túnel SSH.}

\section{Preparación del servidor}\label{sec:desp-preparacion}

Se realiza una sola vez. Los requisitos están en el \cref{cap:requisitos};
el controlador exige Python 3.11 o posterior.

\begin{enumerate}
  \item Como administrador, cree la cuenta de despliegue e instale nginx
    con su servicio global enmascarado
    (\texttt{sudo systemctl mask nginx.service}).
  \item Como usuario de despliegue, desde una copia revisada del
    repositorio, prepare la raíz y valide el resultado:
\begin{verbatim}
loginctl enable-linger "$(id -un)"
python3 -B ops/deploy/host.py --root "$HOME/qadra" --start-databases
systemd-analyze --user verify \
  "$HOME/.config/systemd/user/"qadra-*.service
nginx -t -c "$HOME/qadra/config/nginx.conf"
\end{verbatim}
    El guion verifica las dependencias, crea los directorios base de la
    \cref{tab:desp-raiz}, genera \texttt{settings.json} (lo conserva si ya
    existe), copia los controladores, escribe las unidades y arranca las
    bases. No inicia la aplicación.
  \item Instale y apruebe el lanzador de controladores
    (\cref{sec:desp-controladores}).
\end{enumerate}

La primera activación aplica las migraciones, genera la CA, el certificado
firmante, la TSA y la CRL internas, y publica esa confianza en la base de
datos. No crea ninguna cuenta; las activaciones siguientes conservan claves
y datos.

\section{Construcción del paquete}\label{sec:desp-paquete}

\archivo{scripts/build-release.sh} recibe la versión y el commit, que debe
coincidir con \texttt{HEAD}:

\begin{verbatim}
bash scripts/build-release.sh v1.0.0 "$(git rev-parse HEAD)"
\end{verbatim}

El guion prepara qpdf, verifica la instalación de FFmpeg en
\texttt{/opt/tt-media} (\archivo{docs/media-decoder-setup.md}), compila
backend y frontend con sus archivos de bloqueo y llama a
\archivo{ops/deploy/bundle.py}, que exige un árbol de trabajo sin cambios y
deja en \archivo{output/releases/} el archivo
\texttt{qadra-VERSION-COMMIT.tar.gz} y su \texttt{.sha256}; no sobrescribe
un paquete existente. El paquete reúne \texttt{bin/}
(\texttt{despacho-cli}, \texttt{qadra-ca}, \texttt{ffmpeg} y
\texttt{ffprobe}), \texttt{lib/} (qpdf), \texttt{web/}, \texttt{ca-web/},
\texttt{ca/}, \texttt{migrations/}, \texttt{pki/}, los guiones de
\texttt{ops/} y el manifiesto \texttt{release.json} con versión, commit,
huella de migraciones y SHA-256 de cada archivo. Debe compilarse en una
distribución compatible con la del servidor (\cref{cap:requisitos}).

El flujo \archivo{.github/workflows/deploy.yml} automatiza estos pasos al
publicar una etiqueta de versión (\cref{tab:mant-flujos}).

\section{Activación de una versión}\label{sec:desp-activacion}

Los controladores se invocan siempre a través del lanzador, con un intérprete
aislado y la huella aprobada del inventario de controladores:

\begin{verbatim}
PYTHON -I -B -S RAIZ/controller_launcher.py \
  --root RAIZ --inventory-sha256 SHA256_APROBADO \
  --entrypoint release.py -- --root RAIZ ORDEN
\end{verbatim}

\texttt{PYTHON} es la ruta absoluta del intérprete aprobado y \texttt{ORDEN}
es \texttt{status}, \texttt{activate} o \texttt{rollback}. Cuando el
despliegue lo inicia una etiqueta, \archivo{scripts/deploy-via-ssh.sh}
ejecuta los pasos 1 a 3.

\begin{enumerate}
  \item Verifique el paquete en el equipo de origen con
    \texttt{sha256sum --check} sobre su archivo \texttt{.sha256}.
  \item Copie el paquete a \texttt{RAIZ/incoming/} con \texttt{scp}.
  \item Ejecute la orden de activación, con la ruta del paquete, su
    SHA-256, la versión y el commit:
\begin{verbatim}
activate RAIZ/incoming/PAQUETE.tar.gz SHA256_PAQUETE v1.0.0 COMMIT
\end{verbatim}
  \item Confirme con la orden \texttt{status}, que comprueba la salud de
    la versión activa y devuelve versión, commit y huella de migraciones.
\end{enumerate}

Bajo \texttt{deploy.lock}, \archivo{ops/deploy/release.py} admite el
paquete contra su huella y su manifiesto y lo extrae en
\texttt{releases/}; comprueba versión y huella de migraciones, detiene
\texttt{qadra-web} y \texttt{qadra-api}, toma un respaldo, inicializa la versión, apunta \texttt{current} a ella, arranca
y comprueba su salud. Solo entonces apunta \texttt{previous} a la versión
que estaba activa. Rechaza, antes de detener servicios, una versión menor que la
activa, una versión ya ligada a otro commit y un paquete cuya huella de
migraciones difiera de \texttt{config/schema}.

\importante{Un cambio de esquema no forma parte de la activación
automática: requiere el mantenimiento explícito de
\archivo{docs/database-operations.md}. No edite \texttt{config/schema}
para forzar la activación.}

\section{Reversión}\label{sec:desp-reversion}

Si la activación falla, el controlador restaura y arranca la versión
anterior y devuelve error; si tampoco lo logra, deja detenida la entrada web
(revise \texttt{journalctl --user -u qadra-api}).

Para volver de forma explícita, ejecute la orden \texttt{rollback} con la
misma invocación del lanzador. Exige que exista \texttt{previous} y que
coincida la huella de migraciones, y toma un respaldo como cualquier
activación; al terminar, \texttt{current} y
\texttt{previous} quedan intercambiados. La reversión cambia el código, no
los datos: las escrituras hechas con la versión nueva se conservan. La
restauración de datos se trata en el \cref{cap:operacion}.

\section{Controladores, lanzador y aprobación}\label{sec:desp-controladores}

Los controladores (\texttt{tools/}) se actualizan por un camino separado:
activar una versión no los cambia, ni cambia las unidades. Una
\emph{generación} es el conjunto completo de esos guiones, identificado por
el SHA-256 de su inventario. El operador aprueba esa huella fuera del
servidor y la entrega en cada invocación; nunca se deduce de lo que haya en
el servidor.

El lanzador (\archivo{ops/controller_launcher.py}), la publicación, la
aprobación, el cierre de entradas y el instalador se documentan en los
archivos \archivo{docs/deployment-controller-*.md}. Una generación se
actualiza así:

\begin{enumerate}
  \item Prepare, desde una revisión aprobada, un área privada con lanzador,
    \emph{bootstrap} y candidato, y sus huellas
    (\archivo{docs/deployment-controller-bootstrap.md}).
  \item Suspenda el acceso administrativo externo y deshabilite el
    temporizador de la CRL durante toda la ventana:
    \texttt{systemctl --user disable --now qadra-crl-renew.timer}.
  \item Ejecute el instalador (\texttt{controller\_installation.py}) con
    el lanzador. El estado final es \texttt{installed\_closed}: fuentes
    aprobadas y servicios cerrados. Ante una respuesta incierta se
    reintenta la misma operación.
  \item Entregue la autorización de reapertura: el instalador arranca
    bases, API y web, comprueba la versión y devuelve \texttt{installed}.
  \item Actualice la huella aprobada donde se invoque el lanzador y vuelva
    a generar las unidades del temporizador con la orden
    \texttt{schedule-units} de \texttt{renew\_crl.py} (\cref{sec:op-crl}).
\end{enumerate}

\importante{No use \archivo{ops/deploy/host.py} para actualizar
controladores en una instalación activa: reescribe configuración y unidades,
y las que genera invocan \texttt{tools/runtime.py} sin lanzador ni huella.}
